\documentclass[10pt,a4paper]{amsart}
\usepackage{amsmath,amssymb,amsthm,graphicx}
% For commutative diagrams we use Xy-pic
\usepackage[all]{xy}

% We use 2cell for 2-commutative diagrams.
\xyoption{2cell}
% Theorem environments.
%
\theoremstyle{plain}
\newtheorem{theorem}[subsection]{Theorem}
\newtheorem{proposition}[subsection]{Proposition}
\newtheorem{lemma}[subsection]{Lemma}

\theoremstyle{definition}
\newtheorem{definition}[subsection]{Definition}
\newtheorem{example}[subsection]{Example}
\newtheorem{exercise}[subsection]{Exercise}
\newtheorem{situation}[subsection]{Situation}

\theoremstyle{remark}
\newtheorem{remark}[subsection]{Remark}
\newtheorem{remarks}[subsection]{Remarks}

\numberwithin{equation}{subsection}

% Macros
%
\def\lim{\mathop{\mathrm{lim}}\nolimits}
\def\colim{\mathop{\mathrm{colim}}\nolimits}
\def\Spec{\mathop{\mathrm{Spec}}}
\def\Hom{\mathop{\mathrm{Hom}}\nolimits}
\def\Ext{\mathop{\mathrm{Ext}}\nolimits}
\def\SheafHom{\mathop{\mathcal{H}\!\mathit{om}}\nolimits}
\def\SheafExt{\mathop{\mathcal{E}\!\mathit{xt}}\nolimits}
\def\Sch{\mathit{Sch}}
\def\Mor{\mathop{\mathrm{Mor}}\nolimits}
\def\Ob{\mathop{\mathrm{Ob}}\nolimits}
\def\Sh{\mathop{\mathit{Sh}}\nolimits}
\def\NL{\mathop{N\!L}\nolimits}
\def\CH{\mathop{\mathrm{CH}}\nolimits}
\def\proetale{{pro\text{-}\acute{e}tale}}
\def\etale{{\acute{e}tale}}
\def\QCoh{\mathit{QCoh}}
\def\Ker{\mathop{\mathrm{Ker}}}
\def\Im{\mathop{\mathrm{Im}}}
\def\Coker{\mathop{\mathrm{Coker}}}
\def\Coim{\mathop{\mathrm{Coim}}}

% Boxtimes
%
\DeclareMathSymbol{\boxtimes}{\mathbin}{AMSa}{"02}

%
% Macros for moduli stacks/spaces
%
\def\QCohstack{\mathcal{QC}\!\mathit{oh}}
\def\Cohstack{\mathcal{C}\!\mathit{oh}}
\def\Spacesstack{\mathcal{S}\!\mathit{paces}}
\def\Quotfunctor{\mathrm{Quot}}
\def\Hilbfunctor{\mathrm{Hilb}}
\def\Curvesstack{\mathcal{C}\!\mathit{urves}}
\def\Polarizedstack{\mathcal{P}\!\mathit{olarized}}
\def\Complexesstack{\mathcal{C}\!\mathit{omplexes}}
% \Pic is the operator that assigns to X its picard group, usage \Pic(X)
% \Picardstack_{X/B} denotes the Picard stack of X over B
% \Picardfunctor_{X/B} denotes the Picard functor of X over B
\def\Pic{\mathop{\mathrm{Pic}}\nolimits}
\def\Picardstack{\mathcal{P}\!\mathit{ic}}
\def\Picardfunctor{\mathrm{Pic}}
\def\Deformationcategory{\mathcal{D}\!\mathit{ef}}
\usepackage{cleveref}
\crefname{originallemma}{original lemma}{original lemmas}
\Crefname{originallemma}{Original lemma}{Original lemmas}
\makeatletter
\newlabel{topologies-lemma-at-the-bottom}{{020Z}{source}{}{sourcepointer-1}{}}
\newlabel{topologies-lemma-at-the-bottom@cref}{{[originallemma][0][]020Z}{[1][0][]source}{}{}{}}
\makeatother
\crefname{originallemma}{original lemma}{original lemmas}
\Crefname{originallemma}{Original lemma}{Original lemmas}
\makeatletter
\newlabel{topologies-lemma-at-the-bottom-etale}{{021G}{source}{}{sourcepointer-2}{}}
\newlabel{topologies-lemma-at-the-bottom-etale@cref}{{[originallemma][0][]021G}{[1][0][]source}{}{}{}}
\makeatother
\crefname{originallemma}{original lemma}{original lemmas}
\Crefname{originallemma}{Original lemma}{Original lemmas}
\makeatletter
\newlabel{topologies-lemma-put-in-T}{{020Y}{source}{}{sourcepointer-3}{}}
\newlabel{topologies-lemma-put-in-T@cref}{{[originallemma][0][]020Y}{[1][0][]source}{}{}{}}
\makeatother
\crefname{originallemma}{original lemma}{original lemmas}
\Crefname{originallemma}{Original lemma}{Original lemmas}
\makeatletter
\newlabel{topologies-lemma-morphism-big-small}{{0211}{source}{}{sourcepointer-4}{}}
\newlabel{topologies-lemma-morphism-big-small@cref}{{[originallemma][0][]0211}{[1][0][]source}{}{}{}}
\makeatother
\crefname{originallemma}{original lemma}{original lemmas}
\Crefname{originallemma}{Original lemma}{Original lemmas}
\makeatletter
\newlabel{topologies-lemma-put-in-T-etale}{{021F}{source}{}{sourcepointer-5}{}}
\newlabel{topologies-lemma-put-in-T-etale@cref}{{[originallemma][0][]021F}{[1][0][]source}{}{}{}}
\makeatother
\crefname{originallemma}{original lemma}{original lemmas}
\Crefname{originallemma}{Original lemma}{Original lemmas}
\makeatletter
\newlabel{topologies-lemma-morphism-big-small-etale}{{021I}{source}{}{sourcepointer-6}{}}
\newlabel{topologies-lemma-morphism-big-small-etale@cref}{{[originallemma][0][]021I}{[1][0][]source}{}{}{}}
\makeatother
\crefname{originallemma}{original lemma}{original lemmas}
\Crefname{originallemma}{Original lemma}{Original lemmas}
\makeatletter
\newlabel{lemma-where-sections-are-equal}{{09XM}{source}{}{sourcepointer-7}{}}
\newlabel{lemma-where-sections-are-equal@cref}{{[originallemma][0][]09XM}{[1][0][]source}{}{}{}}
\makeatother
\crefname{originallemma}{original lemma}{original lemmas}
\Crefname{originallemma}{Original lemma}{Original lemmas}
\makeatletter
\newlabel{descent-lemma-open-fpqc-covering}{{03N0}{source}{}{sourcepointer-8}{}}
\newlabel{descent-lemma-open-fpqc-covering@cref}{{[originallemma][0][]03N0}{[1][0][]source}{}{}{}}
\makeatother
\title{The pullback along the projection from the big site to the small site, described by sections and by fpqc sheafiness}
\author{The Stacks Project Authors}
\date{stacks-project:etale-cohomology.tex:local-snapshot-20261011; GNU FDL 1.2 or later}
\begin{document}
\maketitle
\noindent\textit{Editorial note.} All delivered mathematical text is the author's own, copied from the pinned Stacks Project chapter file etale-cohomology.tex; no statement or proof step was rewritten, reordered or removed.\par\medskip
\noindent\textit{Editorial note.} The proof is the author's complete argument for both assertions, including the refinement step in which the largest agreement loci are made to cover.\par\medskip
\noindent\textit{Editorial note.} Every labelled reference occurring inside the counted spans is carried in assembly.external\_references as a non-load-bearing navigation pointer (same-chapter ones at their definition site, cross-chapter ones at the invocation inside the counted span), so no load-bearing mathematics is replaced by a stub. No bibliography citation occurs inside the counted spans.\par\medskip
\noindent\textit{Editorial note.} The checked chapter list for same-name collisions is the set of stacks chapters staged in this repository: chow, cohomology, divisors, etale-cohomology, flat, more-algebra, more-morphisms, morphisms, spaces-more-morphisms, varieties together with this card's chapter etale-cohomology.\par\medskip
\noindent\textit{Editorial note.} The preamble is the chapter's own native preamble excerpt already used for this source: the Xy-pic setup (native preamble lines 14-18) and the native theorem declarations with the mathematical macros (native preamble lines 154-217).\par\medskip
\noindent\textit{Allowed standard prerequisites (editorial).}\begin{itemize}
\item Morphism pi\_S from the big Zariski or etale site of S to the small site. Topologies, the morphisms of sites pi\_S : (Sch/S)\_tau -> S\_tau and their description. Supplies the morphism whose inverse image is described in the statement.
\item Section functor and inverse image along a morphism of schemes. Same chapter, the inverse image f\_\{small\}\textasciicircum{}\{-1\} of a sheaf on S\_tau along f : T -> S and the resulting description of sections of pi\_S\textasciicircum{}\{-1\}F. Supplies the displayed formula and the sheaf axioms used in the proof.
\item fpqc coverings and the sheaf condition for them. Same chapter, the fpqc topology on the big site and the criterion for a presheaf to satisfy the sheaf condition for fpqc coverings. Supplies the second assertion of the statement, which is verified by a gluing argument.
\end{itemize}\medskip
\ifdefined\hypertarget\hypertarget{sourcepointer-1}{}\fi
\noindent\textit{Source reference.} Original lemma 020Z: sources/stacks-etale-cohomology/pinned-original-source.tex, source lines 4965--4965. Navigation pointer invoked at proof line 4965; the counted span contains the pointer only and the cited result is not reproduced inside it, so the pointer is registered as a non-load-bearing navigation stub in the established form for Stacks material. No bibliography citation is involved.\par\medskip
\ifdefined\hypertarget\hypertarget{sourcepointer-2}{}\fi
\noindent\textit{Source reference.} Original lemma 021G: sources/stacks-etale-cohomology/pinned-original-source.tex, source lines 4966--4966. Navigation pointer invoked at proof line 4966; the counted span contains the pointer only and the cited result is not reproduced inside it, so the pointer is registered as a non-load-bearing navigation stub in the established form for Stacks material. No bibliography citation is involved.\par\medskip
\ifdefined\hypertarget\hypertarget{sourcepointer-3}{}\fi
\noindent\textit{Source reference.} Original lemma 020Y: sources/stacks-etale-cohomology/pinned-original-source.tex, source lines 4979--4979. Navigation pointer invoked at proof line 4979; the counted span contains the pointer only and the cited result is not reproduced inside it, so the pointer is registered as a non-load-bearing navigation stub in the established form for Stacks material. No bibliography citation is involved.\par\medskip
\ifdefined\hypertarget\hypertarget{sourcepointer-4}{}\fi
\noindent\textit{Source reference.} Original lemma 0211: sources/stacks-etale-cohomology/pinned-original-source.tex, source lines 4980--4980. Navigation pointer invoked at proof line 4980; the counted span contains the pointer only and the cited result is not reproduced inside it, so the pointer is registered as a non-load-bearing navigation stub in the established form for Stacks material. No bibliography citation is involved.\par\medskip
\ifdefined\hypertarget\hypertarget{sourcepointer-5}{}\fi
\noindent\textit{Source reference.} Original lemma 021F: sources/stacks-etale-cohomology/pinned-original-source.tex, source lines 4981--4981. Navigation pointer invoked at proof line 4981; the counted span contains the pointer only and the cited result is not reproduced inside it, so the pointer is registered as a non-load-bearing navigation stub in the established form for Stacks material. No bibliography citation is involved.\par\medskip
\ifdefined\hypertarget\hypertarget{sourcepointer-6}{}\fi
\noindent\textit{Source reference.} Original lemma 021I: sources/stacks-etale-cohomology/pinned-original-source.tex, source lines 4982--4982. Navigation pointer invoked at proof line 4982; the counted span contains the pointer only and the cited result is not reproduced inside it, so the pointer is registered as a non-load-bearing navigation stub in the established form for Stacks material. No bibliography citation is involved.\par\medskip
\ifdefined\hypertarget\hypertarget{sourcepointer-7}{}\fi
\noindent\textit{Source reference.} Original lemma 09XM: sources/stacks-etale-cohomology/pinned-original-source.tex, source lines 4938--4938. Navigation pointer invoked at proof line 4999; the counted span contains the pointer only and the cited result is not reproduced inside it, so the pointer is registered as a non-load-bearing navigation stub in the established form for Stacks material. No bibliography citation is involved.\par\medskip
\ifdefined\hypertarget\hypertarget{sourcepointer-8}{}\fi
\noindent\textit{Source reference.} Original lemma 03N0: sources/stacks-etale-cohomology/pinned-original-source.tex, source lines 5003--5003. Navigation pointer invoked at proof line 5003; the counted span contains the pointer only and the cited result is not reproduced inside it, so the pointer is registered as a non-load-bearing navigation stub in the established form for Stacks material. No bibliography citation is involved.\par\medskip
\begin{lemma}
\label{lemma-describe-pullback}
Let $S$ be a scheme. Let $\tau \in \{Zariski, \etale\}$. Consider the morphism
$$
\pi_S : (\Sch/S)_\tau \longrightarrow S_\tau
$$
of Topologies, Lemma \ref{topologies-lemma-at-the-bottom} or
\ref{topologies-lemma-at-the-bottom-etale}. Let $\mathcal{F}$ be a sheaf on
$S_\tau$. Then $\pi_S^{-1}\mathcal{F}$ is given by the rule
$$
(\pi_S^{-1}\mathcal{F})(T) = \Gamma(T_\tau, f_{small}^{-1}\mathcal{F})
$$
where $f : T \to S$. Moreover, $\pi_S^{-1}\mathcal{F}$ satisfies the
sheaf condition with respect to fpqc coverings.
\end{lemma}

\begin{proof}
Observe that we have a morphism $i_f : \Sh(T_\tau) \to \Sh(\Sch/S)_\tau)$
such that $\pi_S \circ i_f = f_{small}$ as morphisms
$T_\tau \to S_\tau$, see
Topologies, Lemmas \ref{topologies-lemma-put-in-T},
\ref{topologies-lemma-morphism-big-small},
\ref{topologies-lemma-put-in-T-etale}, and
\ref{topologies-lemma-morphism-big-small-etale}.
Since pullback is transitive we see that
$i_f^{-1} \pi_S^{-1}\mathcal{F} = f_{small}^{-1}\mathcal{F}$ as desired.

\medskip\noindent
Let $\{g_i : T_i \to T\}_{i \in I}$ be an fpqc covering.
The final statement means the following: Given a sheaf $\mathcal{G}$
on $T_\tau$ and given sections
$s_i \in \Gamma(T_i, g_{i, small}^{-1}\mathcal{G})$ whose pullbacks
to $T_i \times_T T_j$ agree, there is a unique section $s$ of $\mathcal{G}$
over $T$ whose pullback to $T_i$ agrees with $s_i$.

\medskip\noindent
Let $V \to T$ be an object of $T_\tau$ and let $t \in \mathcal{G}(V)$.
For every $i$ there is a largest open $W_i \subset T_i \times_T V$
such that the pullbacks of $s_i$ and $t$ agree as sections of the pullback
of $\mathcal{G}$ to $W_i \subset T_i \times_T V$, see
Lemma \ref{lemma-where-sections-are-equal}.
Because $s_i$ and $s_j$ agree over $T_i \times_T T_j$ we find
that $W_i$ and $W_j$ pullback to the same open over
$T_i \times_T T_j \times_T V$. By
Descent, Lemma \ref{descent-lemma-open-fpqc-covering}
we find an open $W \subset V$ whose inverse image to $T_i \times_T V$
recovers $W_i$. 

\medskip\noindent
By construction of $g_{i, small}^{-1}\mathcal{G}$ there exists
a $\tau$-covering $\{T_{ij} \to T_i\}_{j \in J_i}$, for each $j$ an
open immersion or \'etale morphism $V_{ij} \to T$, a section
$t_{ij} \in \mathcal{G}(V_{ij})$, and commutative diagrams
$$
\xymatrix{
T_{ij} \ar[r] \ar[d] & V_{ij} \ar[d] \\
T_i \ar[r] & T
}
$$
such that $s_i|_{T_{ij}}$ is the pullback of $t_{ij}$. In other words,
after replacing the covering $\{T_i \to T\}$ by $\{T_{ij} \to T\}$
we may assume there are factorizations $T_i \to V_i \to T$ with
$V_i \in \Ob(T_\tau)$ and sections $t_i \in \mathcal{G}(V_i)$
pulling back to $s_i$ over $T_i$.
By the result of the previous paragraph we find opens $W_i \subset V_i$
such that $t_i|_{W_i}$ ``agrees with'' every $s_j$ over $T_j \times_T W_i$.
Note that $T_i \to V_i$ factors through $W_i$.
Hence $\{W_i \to T\}$ is a $\tau$-covering and the lemma is proven.
\end{proof}

\end{document}
